\subsection{A diagonal block from conjugacy word lengths}

\begin{claim}{A.31}{Proved}{An algebraic extension of the metric diagonal}
Let $E$ be a countable nonabelian simple group. For $b\ne1$,
let $\rho_b(g)$ denote the least number of conjugates of
$b$ or $b^{-1}$ whose product is $g$, with $\rho_b(1)=0$.
Suppose that
\[
 \text{for every }d\ne1\text{ and every integer }M\geq1,
 \text{ there is }b\ne1\text{ with }\rho_b(d)>M.
 \tag{A.31.1}
\]
Then $E$ has an explicit unsolvable countable system with
identity exponent matrix. For every specified connected CW
classifying space $B$ of a group $H$, every subgroup
$K\leq H$, and every epimorphism $q:K\twoheadrightarrow E$,
this system gives an actual aspherical inclusion
$B\hookrightarrow Y$ using countably many one- and two-cells,
with
\[
 \begin{gathered}
 \dim Y\leq\max\{\dim B,2\},\qquad
 H\hookrightarrow\pi_1Y,\qquad
 H_*(B;\mathbb Z)\xrightarrow{\cong}H_*(Y;\mathbb Z),\\
 \langle\!\langle H\rangle\!\rangle_{\pi_1Y}=\pi_1Y,
 \qquad
 \text{no }\pi_1Y\longrightarrow E\text{ restricts to }q\text{ on }K.
 \end{gathered}
 \tag{A.31.2}
\]
Its relative cellular boundary is the identity in the
displayed cell bases. No map from all of $H$ to $E$ is
required.

There is a new actual two-dimensional acyclic resolution
\[
 t^{\mathrm{alg}}:BG^{\mathrm{alg}}\longrightarrow S^2,
 \qquad |G^{\mathrm{alg}}|\leq2^{\aleph_0},
 \tag{A.31.3}
\]
retaining every conclusion of A.28 and A.30 and excluding
all countable nonabelian simple quotients satisfying
(A.31.1). For each infinite simple quotient $E$ of a
nontrivial countable image of $G^{\mathrm{alg}}$, the function
\[
 N_E(g):=\sup_{b\in E-\{1\}}\rho_b(g)
 \tag{A.31.4}
\]
is a finite integer-valued conjugation-invariant norm and
is unbounded on $E$. For every fixed $d\ne1$, the
conjugacy word norm $\rho_d$ is also unbounded. Thus
$G^{\mathrm{alg}}$ has no countable nonabelian simple
quotient with finite diameter in even one such norm.
Every positive definite
conjugation-invariant length on such an $E$ is discrete.
These quotients also retain all the restrictions proved in
A.28 and A.29, including unimodular closure and
topologizability.
\end{claim}

\begin{proof}
Simplicity makes every $\rho_b(g)$ finite. The function
$\rho_b$ is symmetric, subadditive, positive away from $1$,
and invariant under conjugation; these statements follow
by reversing, concatenating, and conjugating its defining
products. We continue to use $[b,u]=bub^{-1}u^{-1}$.

For $b\ne1$, put
\[
 S_b=\{v[b,u]v^{-1}:u,v\in E\},\qquad
 T_b(u,v)=v[b,u]v^{-1}.
 \tag{A.31.5}
\]
The direct computations in A.30 show that $1\in S_b$,
that $S_b$ is symmetric and invariant under conjugation,
and that it generates $E$. In detail, symmetry follows
from $[b,u]^{-1}=u[b,u^{-1}]u^{-1}$; some commutator
$[b,u]$ is nontrivial because $Z(E)=1$; and the subgroup
generated by $S_b$ is therefore a nontrivial normal
subgroup. Its length estimate now reads
\[
 \rho_b(s)\leq2\qquad(s\in S_b),
 \tag{A.31.6}
\]
because $[b,u]=b(ub^{-1}u^{-1})$.

\paragraph{The finite replacement step.}
Let $R:E^m\to E$ be a nonconstant finite word map, and
let $L\geq1$ be the number of variable occurrences in a
displayed word for $R$, counting inverse occurrences.
For $\mathbf s\in S_b^m$ and $\mathbf u\in E^m$, one has
\[
 \rho_b\bigl(R(\mathbf s\mathbf u)R(\mathbf u)^{-1}\bigr)
 \leq 2L.
 \tag{A.31.7}
\]
Here products of tuples are coordinatewise. To check the
factor order, let $P_0=R(\mathbf u)$ and replace the
variable occurrences one at a time, ending with
$P_L=R(\mathbf s\mathbf u)$. At a positive occurrence,
the replacement of $u_i$ by $s_i u_i$ changes the value by
a conjugate of $s_i$ on the left. At a negative occurrence,
the replacement of $u_i^{-1}$ by $u_i^{-1}s_i^{-1}$
changes the value by a conjugate of $s_i^{-1}$ on the
left. Thus each $P_jP_{j-1}^{-1}$ has $\rho_b$-length
at most two. The exact telescoping product
\[
 P_LP_0^{-1}
 =(P_LP_{L-1}^{-1})(P_{L-1}P_{L-2}^{-1})
       \cdots(P_1P_0^{-1})
\]
proves (A.31.7).

Fix $e\in E$. Choose $\mathbf a$ with
$c:=R(\mathbf a)\ne e$, and choose $\mathbf a'$ with
$R(\mathbf a')\ne c$. By (A.31.1), choose $b\ne1$ such
that
\[
 \rho_b(ec^{-1})>2L.
 \tag{A.31.8}
\]
Since $S_b^m$ generates $E^m$, there is a finite path from
$\mathbf a$ to $\mathbf a'$ whose steps are left
multiplications by elements of $S_b^m$. Choose the first
step at which the value of $R$ changes from $c$. Its
endpoints are $\mathbf u$ and $\mathbf s\mathbf u$, with
\[
 R(\mathbf u)=c\ne R(\mathbf s\mathbf u),
 \qquad\mathbf s\in S_b^m.
 \tag{A.31.9}
\]
For any $\mathbf t\in S_b^m$, equality
$R(\mathbf t\mathbf u)=e$ would make (A.31.7) give
$\rho_b(ec^{-1})\leq2L$, contradicting (A.31.8).
Consequently
\[
 R(S_b^m\mathbf u)\subseteq R(E^m)-\{e\}
\]
contains at least the two distinct values in (A.31.9).
Replacing each old leaf $x_i$ by
$T_b(y_i,z_i)u_i$, with separate new variables for every
pair, gives a finite word whose image is exactly this set.
This proves the finite replacement step without a metric.

\paragraph{The countable system and its actual geometry.}
Enumerate $E$ as $e_0,e_1,\ldots$ and start with
$R_0(x_{\varnothing})=x_{\varnothing}$. Apply the finite
replacement step at level $n$ to the leaves of $R_n$ and
the value $e_n$. On the binary tree of finite words
$s\in\{0,1\}^{<\omega}$ this produces coefficients
$b_n\ne1$ and $u_s\in E$ and the equations
\[
 x_s=T_{b_{|s|}}(x_{s0},x_{s1})u_s.
 \tag{A.31.10}
\]
Every finite substituted root word remains nonconstant;
its image is contained in the preceding image and the
image at level $n+1$ excludes $e_n$. A solution would
place its root value in all these images, although every
element of $E$ is excluded at a finite stage. This is
impossible. Condition (A.31.1) also implies that $E$ is
infinite, since all the finitely many lengths would be
bounded for a finite simple group.

Choose, through $q$, lifts $\beta_n\in K$ and
$\upsilon_s\in K$ of the displayed coefficients, and
based loops in $B$ representing them. Attach a loop
$x_s$ and a two-cell with attaching word
\[
 x_s\upsilon_s^{-1}
   \bigl(x_{s1}[\beta_{|s|},x_{s0}]x_{s1}^{-1}\bigr)^{-1}
 \tag{A.31.11}
\]
for each vertex $s$. The parent has exponent one and both
children have exponent zero, so the relative cellular
boundary sends each two-cell basis vector to the
corresponding $x_s$ basis vector.

Let $Y_n$ contain $B$, the loops at depths at most $n$,
and the two-cells at depths less than $n$. Remove each
internal generator together with its defining two-cell,
in increasing depth order starting at the root. At the
time of its removal, that generator occurs in no other
remaining two-cell: its parent cell has already been
removed, and all coefficients are loops in $B$. The
usual elementary presentation homotopy equivalence fixes
$B$ and the frontier loops and gives
\[
 Y_n\simeq_B B\vee\bigvee_{|s|=n}S^1.
 \tag{A.31.12}
\]
Thus every $Y_n$ is aspherical and $H\to\pi_1Y_n$ is
the inclusion of the indicated free factor. The actual
union $Y=\bigcup_nY_n$ is aspherical because every sphere
of dimension at least two is supported in a finite
subcomplex, hence in one $Y_n$. The same compactness
argument for nullhomotopies proves $H\hookrightarrow
\pi_1Y$. The relative boundary already proves the
homology assertion, and only one- and two-cells were
attached. Killing $H$ makes every coefficient in
(A.31.11) equal to $1$, so every $x_s$ equals $1$ in the
resulting quotient. Finally, a homomorphism to $E$
extending $q$ on $K$ would assign a solution to
(A.31.10). This proves all of (A.31.2).

\paragraph{The resulting resolution and the remaining alternative.}
Rerun the construction of A.28, interleaving its original
three target classes and its one-variable unimodular
adjunctions with all countable targets satisfying
(A.31.1), using (A.31.11) for the new blocks. Groups on
countable underlying sets, countable subgroups of a group
of cardinality at most $2^{\aleph_0}$, and their maps to
countable groups each have at most $2^{\aleph_0}$ choices.
Each block attaches only countably many cells. The same
$\omega_1$ passes as in A.28 therefore keep the total
cardinality at most $2^{\aleph_0}$.

Every map from the final group onto a specified countable
target is detected on a countable subgroup of a single
earlier pass: select one preimage of each target element
and use the uncountable cofinality of $\omega_1$. The
corresponding block in a later pass prohibits that map.
Actual CW unions retain asphericity, homology and normal
generation as in A.28. Starting with its specified A.10
resolution, map each new loop constantly to the basepoint
of $S^2$ and extend each new two-cell over a filling disk;
the attaching image is nullhomotopic since $S^2$ is simply
connected. This extends the original map literally.
The resulting integral homology equivalence to $S^2$ is
acyclic, as verified for the same construction in A.28.
This proves (A.31.3) and all its inherited assertions.

Every target of A.30 satisfies (A.31.1). Indeed, for its
length $\ell$ one has
$\ell(d)\leq\rho_b(d)\ell(b)$; for any prescribed $M$
choose $b\ne1$ with $M\ell(b)<\ell(d)$. Thus the new
construction retains that exclusion as well.

It remains to identify the exact alternative to
(A.31.1). Its failure gives $d\ne1$ and a finite integer
$M$ such that $\rho_b(d)\leq M$ for every $b\ne1$.
For any $g\in E$, choose a finite expression of $g$ as a
product of $k$ conjugates of $d^{\pm1}$. Substituting an
expression for $d$ in conjugates of $b^{\pm1}$ gives
\[
 \rho_b(g)\leq kM\quad\text{for all }b\ne1.
 \tag{A.31.13}
\]
Therefore $N_E(g)$ in (A.31.4) is finite for every $g$.
The supremum is over integers, so it is an integer and is
attained. Taking suprema in the symmetry, subadditivity
and conjugation identities for $\rho_b$ proves that
$N_E$ is a positive definite conjugation-invariant norm.
Boundedness of $N_E$ is exactly uniform simplicity: a
bound supplies the same finite conjugacy word length for
every pair $g,b\ne1$, and such a uniform bound bounds
$N_E$.

For any fixed $d\ne1$, a shortest expression of $g$ in
conjugates of $d^{\pm1}$ and (A.31.4) give the exact comparison
\[
 \rho_d(g)\leq N_E(g)\leq N_E(d)\rho_d(g).
 \tag{A.31.14}
\]
The right inequality follows by replacing each conjugate
of $d^{\pm1}$ by at most $N_E(d)$ conjugates of
$b^{\pm1}$, for each $b\ne1$, and then taking the
supremum. Thus a finite diameter for one $\rho_d$
would bound $N_E$.

The A.28 normal-generation argument supplies an infinite
simple quotient of every nontrivial countable image of
$G^{\mathrm{alg}}$. Any such quotient is excluded from
(A.31.1), and it is also excluded from the uniformly
simple class. Thus its $N_E$ is finite everywhere and
unbounded, and (A.31.14) makes every $\rho_d$ unbounded.
For any other positive definite invariant
length $\ell$, fix $d\ne1$. Every expression of $d$ in
conjugates of $b^{\pm1}$ yields
\[
 0<\ell(d)\leq\rho_b(d)\ell(b)
       \leq N_E(d)\ell(b)\qquad(b\ne1).
 \tag{A.31.15}
\]
Hence $\inf_{b\ne1}\ell(b)\geq\ell(d)/N_E(d)>0$.
This proves discreteness of all these lengths. The
equation-solving and image restrictions pass from A.28,
and A.29 applies to their simple countable quotients.
No boundedness of $N_E$ is inferred from these remaining
properties.
\end{proof}
